% ============================================================================
% docs/informe/Informe-Grupo4.tex -- Informe de implementacion (Proyecto 2,
% entregable 3). CE-1115 Seguridad de la Informacion -- Grupo 4
%
% Un solo archivo: todas las secciones van aqui adentro, cada una con su
% bloque de comentarios y su responsable. Las capturas van en docs/informe/img.
%
% Este archivo reemplaza a docs/Informe_P2.md, docs/P2_Desviaciones_respecto_P1.md
% y docs/Preguntas_abiertas.md: su contenido ya esta aqui (las desviaciones
% completas, como ADR, en el Anexo A).
%
%   Seccion del enunciado (5.1)            Archivo                       Responsable  Paginas
%   0. Datos de entrega                    secciones/00-datos.tex        Meibel       0,5
%   1. Arquitectura desplegada             secciones/01-arquitectura.tex Adriel       1,5
%   2. Controles implementados             secciones/02-controles.tex    Henry        2
%   3. Gestion de secretos                 secciones/03-secretos.tex     Carlos       1
%   4. Seguridad heredada de la plataforma secciones/04-plataforma.tex   Carlos       1
%   5. Desviaciones respecto al diseno     secciones/05-desviaciones.tex Adriel       0,5
%   6. Cobertura de la matriz              secciones/06-cobertura.tex    Ludwin       2
%   7. Evidencia                           secciones/07-evidencia.tex    Meibel       1,5
%                                                                         Total       10
%
% El enunciado pide de 6 a 10 paginas. Se cuentan desde la seccion 1; la
% portada y el indice quedan fuera. Si una seccion se pasa de su
% presupuesto, se recorta ella: no se le quita espacio a otra.
%
% Las secciones de Henry (2) y Ludwin (6) ya traen un borrador pasado desde
% el material que habia en docs/Informe_P2.md: revisarlo, ajustarlo y recortarlo.
%
% Las capturas van en img/. Si falta una, el documento compila igual y
% muestra un recuadro gris con el nombre del archivo.
%
% \pendiente{...} marca en rojo lo que falta. Antes de entregar no debe quedar
% ninguno:   grep -n "pendiente{" Informe-Grupo4.tex
%
% Compilar (desde docs/informe):  pdflatex Informe-Grupo4.tex  (dos veces)
%                       o bien:  latexmk -pdf Informe-Grupo4.tex
% (o en Overleaf: subir la carpeta completa y compilar este archivo)
% ============================================================================
\documentclass[11pt,a4paper]{article}

% ============================================================================
% 1. METADATOS DEL DOCUMENTO
% ============================================================================
\newcommand{\docInstitute}{Tecnol\'ogico de Costa Rica}
\newcommand{\docSchool}{Escuela de Ingenier\'ia en Computadores}
\newcommand{\docCourse}{Seguridad de la Informaci\'on}
\newcommand{\docCode}{CE-1115}
\newcommand{\docKicker}{Proyecto 2}
\newcommand{\docTitle}{Informe de implementaci\'on}
\newcommand{\docSubtitle}{Cliente: Verdeazul Carbono S.A. --- Certificaci\'on de captura de carbono, P\'erez Zeled\'on}
\newcommand{\docTagline}{Del dise\~no a un sistema que corre en la nube}
\newcommand{\docProject}{Implementaci\'on y despliegue del sistema de certificaci\'on\\[0.3cm] de captura de carbono de Verdeazul Carbono S.A.}
\newcommand{\docKeywords}{Cloud Run \textbullet\ Neon \textbullet\ Secret Manager \textbullet\ OPA \textbullet\ Ed25519}
\newcommand{\docGroup}{Grupo 4}
\newcommand{\docProfessor}{MSc. Andr\'es F. Vargas Rivera}
\newcommand{\docSemester}{II Semestre, 2026}
\newcommand{\docDate}{20 de noviembre}
\newcommand{\docPlace}{Cartago, Costa Rica}
\newcommand{\docShortTitle}{P2 --- Informe de implementaci\'on --- \docGroup}

% Datos del sistema, para no repetirlos a mano en cada seccion
\newcommand{\urlSistema}{https://verdeazul-edge-igwgmkeipa-uk.a.run.app}
\newcommand{\urlRepo}{https://github.com/CarlosCL13/CE1115-P2-Grupo4-Verdeazul}

% Integrantes: {Nombre}{Carne}
\newcommand{\docMembers}{%
  \member{Meibel Ceciliano}{2020023333}
  \member{Carlos Contreras}{2021476501}
  \member{Ludwin Ramos}{2021032537}
  \member{Henry Nu\~nez}{2022089224}
  \member{Adriel Chaves}{2021031465}}
\newcommand{\docAuthor}{Meibel Ceciliano, Carlos Contreras, Ludwin Ramos, Henry Nu\~nez, Adriel Chaves}

% ============================================================================
% 2. IDIOMA Y FUENTES
% ============================================================================
\usepackage[T1]{fontenc}
\usepackage[utf8]{inputenc}
\usepackage{lmodern}
\usepackage[spanish,es-tabla,es-nodecimaldot,es-noquoting]{babel}
\usepackage[autostyle,spanish=mexican]{csquotes}
\usepackage{microtype}

% ============================================================================
% 3. MATEMATICA Y UNIDADES
% ============================================================================
\usepackage{amsmath}
\usepackage{amssymb}
\usepackage{amsthm}
\usepackage{mathtools}
\usepackage{siunitx}
\sisetup{per-mode=symbol,range-units=single,output-decimal-marker={,}}

% ============================================================================
% 4. COLORES
% ============================================================================
\usepackage[table]{xcolor}
\definecolor{tecnavy}{RGB}{31,56,100}
\definecolor{tecsteel}{RGB}{46,85,150}
\definecolor{teclight}{RGB}{214,227,243}
\definecolor{tecrule}{RGB}{155,184,216}
\definecolor{boxbg}{RGB}{242,246,252}
\definecolor{labelgray}{RGB}{120,120,120}
\definecolor{codebg}{RGB}{248,248,248}
\definecolor{codekw}{RGB}{0,0,180}
\definecolor{codecmt}{RGB}{0,120,60}
\definecolor{codestr}{RGB}{170,55,55}
% Niveles de clasificacion de activos
\definecolor{nivelpublico}{RGB}{220,240,220}
\definecolor{nivelinterno}{RGB}{214,227,243}
\definecolor{nivelconfidencial}{RGB}{255,236,200}
\definecolor{nivelrestringido}{RGB}{250,215,215}

% ============================================================================
% 5. GRAFICOS Y TIKZ
% ============================================================================
\usepackage{graphicx}
\graphicspath{{img/}{figures/}{./}}
\usepackage{tikz}
\usepackage{pgfplots}
\pgfplotsset{compat=1.18}
\usetikzlibrary{arrows.meta,positioning,shapes.geometric,calc,fit,
                backgrounds,patterns,decorations.pathreplacing,babel}

% ============================================================================
% 6. TABLAS, FLOTANTES Y LISTAS
% ============================================================================
\usepackage{float}
\usepackage{rotating}
\usepackage{array}
\usepackage{booktabs}
\usepackage{multirow}
\usepackage{makecell}
\usepackage{longtable}
\usepackage[justification=centering]{caption}
\usepackage{subcaption}
\usepackage{tabularx}
\usepackage{xltabular}
\usepackage{adjustbox}
\usepackage{pdflscape}
\usepackage{enumitem}

% Tipos de columna: P, C, R con ancho fijo; Y flexible alineada a la izquierda
\newcolumntype{P}[1]{>{\raggedright\arraybackslash}p{#1}}
\newcolumntype{C}[1]{>{\centering\arraybackslash}p{#1}}
\newcolumntype{R}[1]{>{\raggedleft\arraybackslash}p{#1}}
\newcolumntype{Y}{>{\raggedright\arraybackslash}X}
\renewcommand{\arraystretch}{1.15}
\setlength{\tabcolsep}{5pt}
\setlist[itemize]{leftmargin=*,noitemsep,topsep=2pt}
\setlist[enumerate]{leftmargin=*,noitemsep,topsep=2pt}

% Etiquetas de continuacion para tablas que se parten entre paginas
\newcommand{\contnext}[1]{%
  \multicolumn{#1}{r@{}}{\footnotesize\itshape Contin\'ua en la p\'agina siguiente}\\}
\newcommand{\contprev}[1]{%
  \multicolumn{#1}{@{}l}{\footnotesize\itshape
    Tabla \thetable{} -- viene de la p\'agina anterior}\\}

% Encabezado de tabla con fondo institucional
\newcommand{\thd}[1]{\cellcolor{tecnavy}\textcolor{white}{\textbf{#1}}}

% Entornos de definicion y observacion
\theoremstyle{definition}
\newtheorem{definicion}{Definici\'on}[section]
\newtheorem{politica}{Pol\'itica}[section]
\theoremstyle{remark}
\newtheorem*{observacion}{Observaci\'on}

% ============================================================================
% 7. CODIGO FUENTE (LISTINGS)
% ============================================================================
\usepackage{listings}
\renewcommand{\lstlistingname}{C\'odigo}
\renewcommand{\lstlistlistingname}{\'Indice de c\'odigos}
\lstdefinestyle{code}{
  basicstyle=\ttfamily\footnotesize,
  keywordstyle=\color{codekw}\bfseries,
  commentstyle=\color{codecmt}\itshape,
  stringstyle=\color{codestr},
  backgroundcolor=\color{codebg},
  numbers=left,
  numberstyle=\tiny\color{labelgray},
  numbersep=8pt,
  frame=single,
  rulecolor=\color{tecrule},
  breaklines=true,
  breakatwhitespace=false,
  postbreak=\mbox{\textcolor{tecsteel}{$\hookrightarrow$}\space},
  showspaces=false,
  showstringspaces=false,
  tabsize=2,
  captionpos=b,
  columns=flexible,
  upquote=true,
  literate={á}{{\'a}}1 {é}{{\'e}}1 {í}{{\'i}}1 {ó}{{\'o}}1 {ú}{{\'u}}1
           {ñ}{{\~n}}1 {Á}{{\'A}}1 {É}{{\'E}}1 {Í}{{\'I}}1 {Ó}{{\'O}}1
           {Ú}{{\'U}}1 {Ñ}{{\~N}}1 {¿}{{?`}}1 {¡}{{!`}}1
}
\lstset{style=code}

% ============================================================================
% 8. DISENO DE PAGINA, TITULOS Y ENCABEZADOS
% ============================================================================
\usepackage[a4paper,margin=1in,headheight=15pt]{geometry}
\usepackage{setspace}
\setstretch{1.10}
\usepackage{ragged2e}

% Titulos de seccion con el color institucional
\usepackage{titlesec}
\titleformat{\section}
  {\Large\bfseries\color{tecnavy}}{\thesection.}{0.6em}{}
\titleformat{\subsection}
  {\large\bfseries\color{tecsteel}}{\thesubsection.}{0.6em}{}
\titleformat{\subsubsection}
  {\normalsize\bfseries\color{tecsteel}}{\thesubsubsection.}{0.6em}{}

% Indice centrado
\usepackage{tocloft}
\renewcommand{\cfttoctitlefont}{\hfill\Large\bfseries\color{tecnavy}}
\renewcommand{\cftaftertoctitle}{\hfill}

% Encabezado y pie de pagina
\usepackage{fancyhdr}
\usepackage{lastpage}
\pagestyle{fancy}
\fancyhf{}
\fancyhead[L]{\footnotesize\textsc{\docCode}}
\fancyhead[R]{\footnotesize\docShortTitle}
\fancyfoot[C]{\footnotesize P\'agina \thepage\ de \pageref{LastPage}}
\renewcommand{\headrulewidth}{0.4pt}
\renewcommand{\footrulewidth}{0pt}

% Numeracion de figuras, tablas y codigos por seccion
\counterwithin{figure}{section}
\counterwithin{table}{section}
\AtBeginDocument{\counterwithin{lstlisting}{section}}

% ============================================================================
% 9. ENLACES
% ============================================================================
\usepackage{hyperref}
\hypersetup{
  colorlinks=true,
  linkcolor=tecnavy,
  citecolor=tecsteel,
  urlcolor=tecsteel,
  bookmarksnumbered=true,
  pdftitle={\docTitle},
  pdfauthor={\docAuthor},
  pdfsubject={\docCourse\ (\docCode) -- \docGroup},
  pdflang={es}
}
\usepackage[spanish]{cleveref}
% Nombres en espanol coherentes con babel es-tabla: "tabla" (femenino), no
% "cuadro" (masculino), para que concuerden los articulos del texto.
\crefname{table}{tabla}{tablas}
\Crefname{table}{Tabla}{Tablas}
\crefname{figure}{figura}{figuras}
\Crefname{figure}{Figura}{Figuras}
\crefname{section}{sección}{secciones}
\Crefname{section}{Sección}{Secciones}
\crefname{subsection}{sección}{secciones}
\Crefname{subsection}{Sección}{Secciones}

% ============================================================================
% 10. COMANDOS PROPIOS
% ============================================================================
% Cajas de aviso cuando falta un archivo o imagen (el documento compila igual)
\newcommand{\missingfilebox}[1]{%
  \par\medskip\noindent
  \fcolorbox{red!60!black}{red!5}{%
    \parbox{\dimexpr\linewidth-2\fboxsep-2\fboxrule}{%
      \centering\ttfamily\small ARCHIVO FALTANTE: \detokenize{#1}}}%
  \par\medskip}

\newcommand{\missingimagebox}[1]{%
  \begingroup
  \setlength{\fboxsep}{0pt}%
  \fcolorbox{tecrule}{boxbg}{%
    \parbox[c][2.8cm][c]{6cm}{\centering\ttfamily\scriptsize
      IMAGEN NO ENCONTRADA\\[3pt]\detokenize{#1}}}%
  \endgroup}

% \safeimage[opciones]{archivo}: incluye la imagen si existe
\newcommand{\safeimage}[2][width=0.75\linewidth]{%
  \IfFileExists{#2}{\includegraphics[#1]{#2}}{%
    \IfFileExists{img/#2}{\includegraphics[#1]{img/#2}}{%
      \IfFileExists{figures/#2}{\includegraphics[#1]{figures/#2}}{%
        \PackageWarning{tecdoc}{Falta la imagen #2}%
        \missingimagebox{#2}}}}}

% \safelisting[opciones]{archivo}: incluye codigo si existe
\newcommand{\safelisting}[2][]{%
  \IfFileExists{#2}{\lstinputlisting[#1]{#2}}{%
    \PackageWarning{tecdoc}{Falta el codigo #2}%
    \missingfilebox{#2}}}

% \notebox[Titulo]{texto}: caja destacada
\newcommand{\notebox}[2][Nota]{%
  \par\medskip\noindent
  \fcolorbox{tecrule}{boxbg}{%
    \parbox{\dimexpr\linewidth-2\fboxsep-2\fboxrule}{%
      \vspace{2pt}\textbf{\color{tecnavy}#1.}\ #2\vspace{2pt}}}%
  \par\medskip}

\newcommand{\code}[1]{\texttt{#1}}
% \ruta{...}: rutas de archivos, endpoints y nombres con guion bajo. Corta la
% linea en / y _ sin que haya que escapar nada. No usarlo dentro de \caption
% ni de otro comando: ahi va \code{...} con \_ escrito a mano.
\usepackage{xurl}
\DeclareUrlCommand\ruta{\urlstyle{tt}}

% \pendiente{texto}: marca en rojo lo que falta por redactar
\newcommand{\pendiente}[1]{%
  \textcolor{codestr}{\textbf{[PENDIENTE:} \textit{#1}\textbf{]}}}
% \pend: marcador corto para celdas angostas de tablas
\newcommand{\pend}{\textcolor{codestr}{\textbf{[PEND.]}}}

% ============================================================================
% 11. PORTADA
% ============================================================================
% Portada sencilla, centrada, con el mismo formato que las tareas del curso:
% institucion, titulo, curso, nombre del proyecto y datos de entrega.
% Cada integrante va en una linea: Nombre --- Carne.
\newcommand{\member}[2]{#1 \quad--\quad #2\\[4pt]}

\newcommand{\makecover}{%
  \thispagestyle{empty}
  \begingroup
  \setlength{\parindent}{0pt}
  \centering
  \vspace*{\fill}

  {\LARGE\bfseries\color{tecsteel} Instituto \docInstitute\par}
  \vspace{0.6cm}
  {\large \docSchool\par}

  \vspace{2.0cm}
  {\huge\bfseries \docKicker\par}
  \vspace{0.5cm}
  {\Large\bfseries \docTitle\par}
  \vspace{0.5cm}
  {\large\itshape \docCode\ \docCourse\par}

  \vspace{1.6cm}
  {\large\bfseries \docProject\par}

  \vspace{1.6cm}
  {\large Estudiantes:\par}
  \vspace{0.3cm}
  {\large \docMembers\par}

  \vspace{1.5cm}
  {\large Profesor: \docProfessor\par}
  \vspace{0.5cm}
  {\large Fecha de entrega: \docDate\par}
  \vspace{0.5cm}
  {\large \docSemester\par}
  \vspace*{\fill}

  \endgroup
  \newpage}

% ============================================================================
% 12. CUERPO DEL DOCUMENTO
% ============================================================================
\begin{document}

\makecover

\tableofcontents
\clearpage

% ============================================================================
% 0. DATOS DE ENTREGA                                   Responsable: Meibel
% Presupuesto: media pagina. Fuente: docs/Informe_P2.md, seccion 0.
% ============================================================================
\section*{Datos de entrega}
\addcontentsline{toc}{section}{Datos de entrega}

\begin{center}
\begin{tabular}{@{}P{0.28\linewidth}P{0.66\linewidth}@{}}
\toprule
\thd{Elemento} & \thd{Valor}\\
\midrule
Sistema desplegado & \url{\urlSistema}\\
Repositorio        & \url{\urlRepo}\\
Google Cloud       & Proyecto \ruta{verdeazul-g4}, regi\'on \ruta{us-east4}\\
Neon               & \ruta{verdeazul-app} (principal) y \ruta{verdeazul-bitacora} (bit\'acora), AWS \ruta{us-east-1}\\
Versi\'on entregada & \pendiente{etiqueta de las im\'agenes desplegadas el d\'ia de la entrega (los 7 caracteres del commit)}\\
\bottomrule
\end{tabular}
\end{center}

\paragraph{Credenciales de prueba.} Hay una cuenta por rol: \ruta{dg},
\ruta{dc}, \ruta{pa}, \ruta{tc1}, \ruta{tc2}, \ruta{vf1}, \ruta{vf2},
\ruta{ae}, \ruta{df1}, \ruta{df2} y \ruta{ta}. Las de \ruta{dg}, \ruta{dc},
\ruta{pa}, \ruta{vf1} y \ruta{vf2} piden adem\'as un c\'odigo TOTP.
\pendiente{decidir con el grupo c\'omo se le entregan las contrase\~nas y las
l\'ineas otpauth al profesor: por un canal aparte o en un anexo de este PDF.
Nunca en el repositorio.}

\clearpage

% ============================================================================
% 1. ARQUITECTURA DESPLEGADA                            Responsable: Adriel
% Presupuesto: 1,5 paginas. Fuente: docs/Informe_P2.md, seccion 1.
%
% Pide el enunciado: "Diagrama de lo que realmente construyeron:
% componentes, imagenes y como se comunican."
%
% Las etiquetas F1..F6 son las fronteras de confianza de la Parte 1, tal como
% las define la tabla 4.x del documento de arquitectura: F3 backend->base,
% F4 backend->motor, F5 navegador<->proveedor de identidad, F6 backend y
% transmisor->bitacora y gestor de secretos.
% ============================================================================
\section{Arquitectura desplegada}
\label{sec:arquitectura}

El sistema corre en Google Cloud Run, en la regi\'on \ruta{us-east4}, como
tres servicios construidos a partir de cuatro im\'agenes de Docker publicadas en
Artifact Registry con la etiqueta de su commit. Los datos viven en dos
proyectos de Neon y la copia de retenci\'on, en Cloud Storage. El \'unico
servicio que responde a internet es el borde; el backend y el transmisor
rechazan cualquier invocaci\'on que no traiga una identidad de Google
autorizada (\cref{fig:arquitectura}).

\begin{figure}[H]
\centering
\resizebox{0.92\linewidth}{!}{%
\begin{tikzpicture}[
  font=\small,
  >={Stealth[length=2.2mm]},
  caja/.style={draw=tecsteel, fill=boxbg, rounded corners=2pt, align=center,
               minimum height=12mm, text width=31mm, inner sep=4pt},
  datos/.style={draw=tecnavy, fill=teclight, cylinder, shape border rotate=90,
               aspect=0.15, align=center, text width=30mm, minimum height=14mm},
  etq/.style={font=\footnotesize, fill=white, inner sep=1.5pt, align=center},
  flecha/.style={->, thick, color=tecnavy},
  punteada/.style={->, dashed, thick, color=labelgray}
]
  % Fuera de Google
  \node[caja, text width=18mm] (nav) at (0,0) {Navegador};

  % Cloud Run
  \node[caja] (edge) at (4.6,0) {\textbf{verdeazul-edge}\\imagen \texttt{edge}\\p\'ublico};
  \node[caja] (api)  at (10.2,0.9) {\textbf{verdeazul-api}\\imagen \texttt{api} (PEP)};
  \node[caja] (opa)  at (10.2,-1.0) {OPA (PDP)\\imagen \texttt{policies}};
  \node[caja, text width=37mm] (tra)  at (10.2,-4.6) {\textbf{verdeazul-transmisor}\\imagen \texttt{transmisor}};
  \begin{scope}[on background layer]
    \node[draw=tecsteel, dashed, rounded corners, fit=(api)(opa), inner sep=5pt] (apibox) {};
    \node[font=\footnotesize, color=tecsteel, anchor=south east] at (apibox.north east) {servicio cerrado};
    \node[draw=tecnavy, thick, rounded corners, fit=(edge)(apibox)(tra), inner sep=12pt] (cr) {};
    \node[font=\footnotesize\bfseries, color=tecnavy, anchor=south west] at (cr.north west) {Cloud Run, us-east4};
  \end{scope}

  % Servicios de Google fuera de Cloud Run
  \node[caja] (sm)  at (10.2,4.2) {Secret Manager};
  \node[caja] (sch) at (4.6,-4.6) {Cloud Scheduler};

  % Datos
  \node[datos] (neon) at (16.6,0.9) {Neon\\\texttt{verdeazul-app}};
  \node[datos] (bit)  at (16.6,-2.2) {Neon\\\texttt{verdeazul-bitacora}};
  \node[datos] (gcs)  at (16.6,-5.2) {Cloud Storage\\retenci\'on 10 a\~nos};

  % Flujos
  \draw[flecha] (nav) -- node[etq, above=1pt]{HTTPS\\(F1)} (edge);
  \draw[flecha] (edge.east) -- node[etq, above=2pt]{HTTPS + token\\de Google (F2)} (api.west);
  \draw[flecha] (api) -- node[etq, right=2pt]{127.0.0.1 (F4)} (opa);
  \draw[flecha] (api.east) -- node[etq, above=2pt]{\texttt{app\_rw} (F3)} (neon.west);
  \draw[flecha] (api.south east) -- node[etq, pos=0.62]{\texttt{log\_reader} (F6)} (bit.north west);
  \draw[flecha] (tra.north east) -- node[etq, pos=0.3]{\texttt{outbox\_rw} (F3)} (neon.south west);
  \draw[flecha] (tra.east) -- node[etq, pos=0.55]{\texttt{log\_writer} (F6)} (bit.south west);
  \draw[flecha] (tra.south east) -- node[etq, below=2pt, pos=0.5]{volcado semanal} (gcs.west);
  \draw[flecha] (sch) -- node[etq, above=2pt]{OIDC} (tra);
  \draw[punteada] (sm) -- node[etq, right=2pt]{\texttt{-{}-set-secrets} (F6)} (api);
  \draw[punteada] (sm.west) -| ([xshift=-9mm]tra.north west) -- (tra.north west);
\end{tikzpicture}}
\caption{Lo que corre en producci\'on. Las flechas llenas son llamadas; las
punteadas, secretos inyectados al arrancar. F1 a F6 son las fronteras de
confianza de la Parte~1; F5 ya no aparece porque no hay proveedor de identidad
aparte (D-05).}
\label{fig:arquitectura}
\end{figure}

% Tabla: componente de la Parte 1 -> donde quedo.
\begin{table}[H]
\centering\footnotesize
\caption{Componentes de la Parte~1 y d\'onde quedaron.}
\label{tab:componentes}
\begin{tabularx}{\linewidth}{@{}P{3.4cm}YP{1.9cm}@{}}
\toprule
\thd{Componente (P1)} & \thd{Aqu\'i} & \thd{Imagen}\\
\midrule
Proxy inverso e interfaz & \ruta{services/edge} y el TLS de Cloud Run; el \'unico p\'ublico (D-04, D-23) & \ruta{edge}\\
Backend (PEP) & \ruta{services/api}, solo invocable por la cuenta del borde & \ruta{api}\\
Motor de pol\'iticas (PDP) & \ruta{services/policies}, contenedor adjunto del backend (D-06) & \ruta{policies}\\
Proveedor de identidad & Dentro del backend: Argon2id, sesi\'on y TOTP (D-05) & \ruta{api}\\
Transmisor y trabajos & \ruta{services/transmisor}: transmitir, vencer, cadena y volcado (D-11) & \ruta{transmisor}\\
Base de datos (PIP) & Neon \ruta{verdeazul-app}, con PostGIS y un rol por servicio & ---\\
Bit\'acora & Neon \ruta{verdeazul-bitacora}: otra instancia (F6) & ---\\
Gestor de secretos & Secret Manager, una cuenta de servicio por componente (D-03) & ---\\
Copia de retenci\'on & Cloud Storage \ruta{verdeazul-g4-volcados}, 10 a\~nos (D-07) & ---\\
\bottomrule
\end{tabularx}
\end{table}

\begin{table}[H]
\centering\footnotesize
\caption{Las fronteras de confianza de la Parte~1 y c\'omo se hacen cumplir en lo desplegado.}
\label{tab:fronteras}
\begin{tabularx}{\linewidth}{@{}P{3.3cm}Y@{}}
\toprule
\thd{Frontera (P1)} & \thd{C\'omo se hace cumplir}\\
\midrule
\textbf{F1} Internet $\rightarrow$ borde &
TLS de Cloud Run (D-04); en el borde, HSTS, CSP y las dem\'as cabeceras,
cookie \ruta{__Host-} con \code{HttpOnly}, \code{Secure} y
\code{SameSite=Strict}, rechazo de formularios de otro origen y tama\~no
m\'aximo de lo que se sube.\\
\addlinespace
\textbf{F2} Borde $\rightarrow$ backend &
El backend no acepta invocaciones an\'onimas: solo la cuenta
\ruta{verdeazul-edge} puede llamarlo. Adentro valida la sesi\'on (opaca y
guardada en la base), la forma de cada petici\'on y el token anti-CSRF que
comprueba el borde.\\
\addlinespace
\textbf{F3} Backend $\rightarrow$ base &
TLS con \code{sslmode=verify-full}, consultas parametrizadas y el rol
\ruta{app_rw}: no es due\~no de las tablas, no tiene DELETE y solo puede
actualizar las columnas que cada flujo necesita. PostGIS valida el pol\'igono
de cada finca.\\
\addlinespace
\textbf{F4} Backend $\rightarrow$ motor &
El motor escucha solo en \ruta{127.0.0.1} dentro del mismo servicio; los
atributos de cada consulta salen de la base, nunca del cliente; si el motor
no responde, se niega todo. Pas\'o de frontera de red a frontera de proceso
(D-06).\\
\addlinespace
\textbf{F5} Navegador $\leftrightarrow$ identidad &
Ya no existe como frontera aparte: sin Keycloak, el inicio de sesi\'on cruza
F1 y F2 como cualquier otra petici\'on, y el segundo factor y la
reautenticaci\'on los exige el backend (D-05).\\
\addlinespace
\textbf{F6} Backend y transmisor $\rightarrow$ bit\'acora y secretos &
La bit\'acora es otra instancia de Neon: el transmisor escribe con
\ruta{log_writer} (solo INSERT) y la API lee con \ruta{log_reader} (solo
SELECT). Cada evento entra a la bandeja de salida en la misma transacci\'on
que la operaci\'on y se sella con un HMAC encadenado. Los secretos llegan de
Secret Manager, a una cuenta de servicio por componente.\\
\bottomrule
\end{tabularx}
\end{table}

\paragraph{Una petici\'on, de punta a punta.} El borde recibe la cookie y
llama al backend con el token de sesi\'on en una cabecera y su identidad de
Google en otra. El backend carga la sesi\'on y las zonas vigentes de la
persona, lee de la base el objeto real que se pide y le pregunta al motor con
el sujeto, la operaci\'on y los atributos. Si el motor dice que s\'i, ejecuta
y anota el evento en la bandeja de salida dentro de la misma transacci\'on; si
dice que no, anota \ruta{ACCESO_DENEGADO} con el motivo y responde 403. Cada
30~minutos el transmisor copia la bandeja a la bit\'acora y encadena cada
evento con HMAC; cada lunes copia la bit\'acora entera a Cloud Storage.

% ============================================================================
% 2. CONTROLES IMPLEMENTADOS                            Responsable: Henry
% Presupuesto: 2 paginas. Fuente: docs/Informe_P2.md, seccion 2.
%
% Pide el enunciado: "Por cada control obligatorio: donde esta en el codigo
% y como se verifica que funciona." Los obligatorios son los seis de 2.3.
%
% La fila 2.3.6 (firma) ya esta escrita: es de Meibel y esta probada en
% produccion. Las otras cinco salen de la tabla de Informe_P2.md.
% Ojo con 2.3.4: Informe_P2.md dice que falta HSTS. Confirmar antes de
% entregar si se agrego o declararlo como pendiente.
% ============================================================================
\section{Controles implementados}
\label{sec:controles}

% Borrador pasado desde docs/Informe_P2.md (seccion 2). Henry: revisar y
% recortar. HSTS ya esta: lo agrego la rama fix/hsts-opa-y-readmes.
Los seis controles obligatorios de la secci\'on~2.3 del enunciado est\'an
implementados en el servidor, y cada uno tiene pruebas autom\'aticas que corren
en cada \emph{pull request} (\cref{tab:controles}).

\begin{xltabular}{\linewidth}{@{}P{2.7cm}YY@{}}
\caption{Controles obligatorios: d\'onde est\'an y c\'omo se verifican.}
\label{tab:controles}\\
\toprule
\thd{Control} & \thd{D\'onde est\'a} & \thd{C\'omo se verifica}\\
\midrule
\endfirsthead
\contprev{3}
\toprule
\thd{Control} & \thd{D\'onde est\'a} & \thd{C\'omo se verifica}\\
\midrule
\endhead
\midrule
\contnext{3}
\endfoot
\bottomrule
\endlastfoot
\textbf{2.3.1} Contrase\~nas con hash lento &
  \ruta{services/api/app/auth/contrasenas.py}: Argon2id con m~=~19~MiB,
  t~=~2, p~=~1 y sal por contrase\~na; m\'inimo de 15 caracteres, lista de
  prohibidas y consulta a filtraciones conocidas (CO-02). &
  \ruta{tests/identidad/test_unitarias.py}: el hash empieza con
  \ruta{$argon2id$v=19$m=19456,t=2,p=1$} y la sal cambia en cada hash. \\
\textbf{2.3.2} Control de acceso en el servidor &
  Cada endpoint llama a \code{autorizar()}
  (\ruta{services/api/app/auth/dependencias.py}) antes de tocar la base; el
  motor en Rego (\ruta{services/policies/rego/}) decide y niega por omisi\'on,
  tambi\'en si no responde. &
  977 pruebas de la matriz contra OPA real (\ruta{tests/matriz}) y pruebas
  por endpoint (\cref{sec:cobertura}). En producci\'on:
  \cref{fig:ev-denegado}. \\
\textbf{2.3.3} Secretos fuera del c\'odigo &
  Secret Manager, inyectados con \ruta{--set-secrets}
  (\ruta{.github/workflows/desplegar.yml}); \ruta{.env} en \ruta{.gitignore};
  despliegue con federaci\'on de identidad, sin archivo de credenciales
  (\cref{sec:secretos}). &
  Job \emph{Buscar credenciales versionadas}: gitleaks sobre todo el
  historial en cada \emph{pull request}. En producci\'on:
  \cref{fig:ev-secretos}. \\
\textbf{2.3.4} TLS en todo el tr\'afico &
  Cloud Run solo sirve HTTPS en \ruta{*.run.app}; el borde agrega
  \code{Strict-Transport-Security}; las conexiones a Neon exigen
  \code{sslmode=verify-full} (\ruta{config.py} rechaza otra cosa); el borde
  llama al backend por HTTPS. &
  \cref{fig:ev-https}; el contenedor no arranca si una cadena de conexi\'on no
  pide \code{verify-full}. \\
\textbf{2.3.5} Validaci\'on de entradas &
  Consultas parametrizadas en todos los m\'odulos; filtros con forma cerrada
  (patrones en \code{Query}, modelos de Pydantic) y fragmentos de SQL fijos. &
  Pruebas de 422 con entradas como \code{X' OR '1'='1} y
  \code{1; DROP TABLE eventos} (\ruta{tests/bitacora}), y en los modelos de
  enmiendas, asignaciones y consentimientos. \\
\textbf{2.3.6} Firma digital de certificados &
  \ruta{certificados/firma.py}: Ed25519 con la llave privada de Secret
  Manager (\ruta{llave-firma-certificados}). Se firma la forma can\'onica del
  contenido con un prefijo de dominio; la clave p\'ublica y su huella viajan en
  el documento. \ruta{verificacion.py} comprueba sin tocar la base y compara la
  huella con la publicada de Verdeazul. La base no deja modificar lo firmado. &
  22 pruebas sin base en \ruta{tests/certificados/test_firma.py} (v\'alido,
  alterado, firmado con otra llave, campos sin firmar) y 41 de punta a punta
  contra PostGIS y OPA. En producci\'on: \cref{fig:ev-valido}.
  Huella publicada: \ruta{SHA256:203b7265...df811d}. \\
\end{xltabular}

% Opcional, si cabe en el presupuesto: controles de P1 que tambien quedaron
% (TOTP, bloqueo por intentos, datos cifrados, separacion de deberes,
% revocar el mismo dia, contenedores sin root). Ver la lista en
% Informe_P2.md, seccion 2. Una linea por control, en una lista.
\paragraph{Controles de la Parte~1 que tambi\'en quedaron.}
\begin{itemize}
  \item Segundo factor TOTP para DG, DC, PA y VF (D-13), y reautenticaci\'on
    antes de aprobar, emitir, anular, abrir y cerrar per\'iodos, cerrar
    enmiendas, cambiar cuentas, asignar zonas y autorizar a terceros.
  \item Bloqueo por intentos (RC-01): espera creciente desde el quinto y
    30~minutos al d\'ecimo.
  \item Datos personales cifrados con AES-256-GCM (LG-05) y b\'usqueda por
    c\'edula con HMAC con llave.
  \item Separaci\'on de deberes din\'amica: C3, C4, C5 (tres personas por
    enmienda) y C10, reforzadas en la base.
  \item Revocar el mismo d\'ia (CU-11): las zonas se leen en cada petici\'on.
  \item Bit\'acora en otra instancia, con cadena HMAC, comprobaci\'on mensual y
    copia semanal con retenci\'on de diez a\~nos.
  \item Contenedores sin root e im\'agenes etiquetadas con su commit, nunca
    \ruta{latest}.
\end{itemize}

% ============================================================================
% 3. GESTION DE SECRETOS                                Responsable: Carlos
% Presupuesto: 1 pagina. Fuente: docs/Informe_P2.md, seccion 3.
%
% Pide el enunciado: "Que secretos existen, donde viven y como llegan al
% contenedor." Nunca poner un valor real, ni parte de uno.
% ============================================================================
\section{Gesti\'on de secretos}
\label{sec:secretos}

El sistema usa ocho secretos. Viven en Secret Manager, en el proyecto
\ruta{verdeazul-g4}, y ninguno est\'a en el repositorio: en cada pull request,
el job \emph{Buscar credenciales versionadas} corre gitleaks sobre todo el
historial de git, que conserva lo que se borra (\cref{tab:secretos}).

\begin{table}[H]
\centering\small
\caption{Secretos en Secret Manager y qui\'en los lee.}
\label{tab:secretos}
\begin{tabularx}{\linewidth}{@{}P{5.6cm}YP{2.4cm}@{}}
\toprule
\thd{Secreto} & \thd{Qu\'e es} & \thd{Lo lee}\\
\midrule
\ruta{verdeazul-db-url} & Cadena de \ruta{app_rw}, base principal & backend\\
\ruta{verdeazul-bitacora-lectura-url} & Cadena de \ruta{log_reader}, bit\'acora & backend y transmisor\\
\ruta{llave-busqueda-cedula} & HMAC para buscar por c\'edula & backend\\
\ruta{llave-cifrado-datos} & AES-256-GCM de los datos personales & backend\\
\ruta{llave-firma-certificados} & Llave privada Ed25519 & backend\\
\ruta{verdeazul-outbox-url} & Cadena de \ruta{outbox_rw}, bandeja de salida & transmisor\\
\ruta{verdeazul-bitacora-escritura-url} & Cadena de \ruta{log_writer}, bit\'acora & transmisor\\
\ruta{llave-cadena-bitacora} & HMAC de la cadena de la bit\'acora & transmisor\\
\bottomrule
\end{tabularx}
\end{table}

\subsection*{C\'omo llegan al contenedor}

El despliegue los conecta con \ruta{--set-secrets}, en la forma
\ruta{VARIABLE=secreto:latest}. La revisi\'on de Cloud Run guarda solo la
referencia al secreto y lo inyecta como variable de entorno al arrancar el
contenedor: no est\'an en el Dockerfile ni en la imagen. El
\ruta{config.py} de cada servicio los envuelve en una clase \code{Secreto} que
se imprime como \code{<secreto oculto>}, as\'i que un registro o una traza de
error no muestran el valor. Si en la nube falta uno, o una cadena de
conexi\'on no pide \code{sslmode=verify-full}, el contenedor no arranca y Cloud
Run sigue sirviendo la revisi\'on anterior.

\subsection*{Qui\'en puede leer cu\'al}

Cada servicio corre con su propia cuenta de servicio, con el rol
\ruta{roles/secretmanager.secretAccessor} solo sobre sus secretos: la del
backend lee cinco, la del transmisor cuatro, y la del borde, el \'unico
servicio abierto a internet, ninguno. Sin cuentas propias, Cloud Run usar\'ia
la cuenta por omisi\'on de Compute, que es Editor de todo el proyecto. El
despliegue tampoco usa una llave: GitHub Actions se autentica con federaci\'on
de identidad (\emph{Workload Identity Federation}), as\'i que no existe un
archivo JSON de cuenta de servicio, ni en el repositorio ni en los secretos de
GitHub.

\subsection*{Riesgos que quedan}

Los registra la desviaci\'on D-03. Las credenciales de base de datos son
est\'aticas: si una se filtra, sirve hasta que alguien la cambie, mientras que
OpenBao las renovaba cada 24 horas. Las llaves del HMAC, del cifrado y de la
firma est\'an en la memoria del proceso que las usa. Y las personas con rol de
administrador del proyecto de Google Cloud pueden leer los ocho secretos
(RR-02 de la Parte~1).

% ============================================================================
% 4. SEGURIDAD HEREDADA DE LA PLATAFORMA                Responsable: Carlos
% Presupuesto: 1 pagina. Fuente: docs/Informe_P2.md, seccion 4.
%
% Pide el enunciado (5.2), para Neon, al menos:
%   - que protege por defecto (cifrado en reposo y en transito, aislamiento,
%     respaldos y recuperacion a un punto, autenticacion de la conexion);
%   - que configuraron (modo de conexion, IP de origen, usuario de la app y
%     sus permisos, ramas o entornos separados);
%   - por que esa configuracion y que alternativa descartaron;
%   - que queda sin cubrir, nombrando las amenazas de la Parte 1.
% Y la misma linea para Cloud Run. "Media pagina basta": se evalua trazar
% la frontera, no la extension. Nunca escribir "Neon es seguro" sin decir
% contra que.
% ============================================================================
\section{Seguridad heredada de la plataforma}
\label{sec:plataforma}

Neon y Cloud Run protegen por su cuenta el transporte, el almacenamiento y la
infraestructura. El c\'odigo, los permisos y los secretos quedaron de nuestro
lado, y para cada plataforma se dice contra qu\'e amenazas de la Parte~1
protege cada lado.

\subsection{Neon}

\begin{table}[H]
\centering\small
\caption{Neon: lo que pone la plataforma y lo que pusimos nosotros.}
\label{tab:neon}
\begin{tabularx}{\linewidth}{@{}P{3.1cm}Y@{}}
\toprule
\thd{Aspecto} & \thd{Detalle}\\
\midrule
Lo que protege por defecto &
Cifrado del almacenamiento en reposo; TLS obligatorio y autenticaci\'on con
SCRAM en toda conexi\'on; cada proyecto con su propio c\'omputo y
almacenamiento; recuperaci\'on a cualquier punto de las \'ultimas seis horas;
el respaldo del almacenamiento lo gestiona el proveedor.\\
\addlinespace
Lo que configuramos &
\textbf{Conexi\'on:} \code{sslmode=verify-full}, exigido en el c\'odigo, que
adem\'as de cifrar verifica el certificado del servidor.
\textbf{Entornos:} dos proyectos, \ruta{verdeazul-app} y
\ruta{verdeazul-bitacora}, para que la bit\'acora sea otra instancia que la
aplicaci\'on no alcanza (F6).
\textbf{Usuario de la aplicaci\'on:} \ruta{app_rw}, sin DELETE, sin cambios
de esquema y con UPDATE solo en las columnas que cada flujo necesita. Cada
servicio tiene su propio rol, y el due\~no, \ruta{neondb_owner}, no lo usa
ninguno.
\textbf{Origen:} sin restringir.\\
\addlinespace
Por qu\'e, y qu\'e se descart\'o &
Dos ramas de un mismo proyecto comparten el usuario due\~no y la cuota; sirven
para separar entornos de trabajo, y la bit\'acora necesitaba otra instancia con
sus propias credenciales. La restricci\'on por direcci\'on de origen es
de plan pago, y Cloud Run no sale a internet por una direcci\'on fija sin un
conector de pago (D-08).\\
\addlinespace
Lo que no cubre (nuestro) &
Una consulta mal construida, y por eso todas son parametrizadas; un permiso de
m\'as, y por eso los permisos por columna; una credencial filtrada, que sin
restricci\'on por origen sirve desde cualquier lugar hasta que se rote (D-03,
D-08); los diez a\~nos de retenci\'on que pide la Parte~1, con una ventana de
seis horas (D-07); y un administrador de la cuenta de Neon, que puede borrar
los proyectos (RR-02).\\
\bottomrule
\end{tabularx}
\end{table}

De las amenazas de la Parte~1, la plataforma cubre AE-03, interceptar la
conexi\'on de un t\'ecnico en una red p\'ublica, con TLS en toda conexi\'on.
Estas quedaron de nuestro lado:

\begin{itemize}
  \item \textbf{AI-02}, modificar una medici\'on aprobada sin dejar rastro:
    \ruta{app_rw} no tiene UPDATE sobre la biomasa, y corregirla exige una
    enmienda con tres personas.
  \item \textbf{AI-07}, alterar la bit\'acora: otra instancia donde solo se
    inserta, un sello HMAC encadenado y una copia semanal que no se puede borrar.
  \item \textbf{AE-01}, cambiar el n\'umero de finca en una petici\'on (IDOR): el
    motor decide con los datos del objeto le\'idos de la base, y el n\'umero que
    manda el cliente solo sirve para buscarlo.
  \item \textbf{AE-02}, contrase\~nas filtradas: Argon2id, lista de prohibidas,
    segundo factor y bloqueo por intentos.
\end{itemize}

\subsection{Cloud Run: responsabilidad compartida}

\begin{table}[H]
\centering\small
\caption{Cloud Run: d\'onde est\'a la l\'inea.}
\label{tab:cloudrun}
\begin{tabularx}{\linewidth}{@{}YY@{}}
\toprule
\thd{Google} & \thd{Nosotros}\\
\midrule
La infraestructura y el aislamiento entre contenedores &
El c\'odigo y sus dependencias, vigiladas por Dependabot\\
El certificado y la terminaci\'on del TLS &
Las cabeceras de seguridad del borde: CSP, HSTS y las dem\'as (D-04)\\
La protecci\'on del frente de red y el escalado &
Qu\'e es p\'ublico: solo el borde. El backend y el transmisor exigen una
identidad de Google (F2)\\
Ejecutar las im\'agenes que se le indican &
Im\'agenes sin root y con la etiqueta de su commit; cuentas de servicio,
permisos y secretos (\cref{sec:secretos})\\
\bottomrule
\end{tabularx}
\end{table}

Lo que P1 le ped\'ia a nginx y Cloud Run no hace: limitar el TLS a la versi\'on
1.3, porque acepta tambi\'en la 1.2, y un l\'imite de tasa general (D-04).

% ============================================================================
% 5. DESVIACIONES RESPECTO AL DISENO                     Responsable: Adriel
% Presupuesto: 0,5 paginas. Fuente: docs/P2_Desviaciones_respecto_P1.md.
%
% Pide el enunciado: "Que cambio y por que", y en 4.1: "que cambiaron, por
% que, y que implicacion de seguridad tuvo". Son 25 (D-01 a D-25): aqui va
% la tabla resumen y un parrafo de las de mas peso; cada una completa, como
% registro de decision de arquitectura (ADR), esta en el Anexo A.
% ============================================================================
\section{Desviaciones respecto al dise\~no}
\label{sec:desviaciones}

El modelo de control de acceso no cambi\'o: los ocho roles, los once objetos,
las once operaciones y las condiciones C1 a C11 se implementaron como se
dise\~naron. Lo que cambi\'o es con qu\'e se construye. Son 25 desviaciones
(\cref{tab:desviaciones}); cada una est\'a escrita completa en el
Anexo~\ref{anx:adr} como un registro de decisi\'on de arquitectura, con su contexto,
lo que se decidi\'o y sus consecuencias positivas y negativas.

{\footnotesize
\setlength{\tabcolsep}{4pt}\renewcommand{\arraystretch}{1.05}
\begin{xltabular}{\linewidth}{@{}P{0.9cm}YP{2.4cm}@{}}
\caption{Desviaciones respecto a la Parte~1 y su implicaci\'on de seguridad.}
\label{tab:desviaciones}\\
\toprule
\thd{} & \thd{Qu\'e cambi\'o y qu\'e implica} & \thd{Origen}\\
\midrule
\endfirsthead
\contprev{3}
\toprule
\thd{} & \thd{Qu\'e cambi\'o y qu\'e implica} & \thd{Origen}\\
\midrule
\endhead
\midrule
\contnext{3}
\endfoot
\bottomrule
\endlastfoot
D-01 & Alojamiento en servicios propietarios. Crece RR-02 (el proveedor custodia secretos y datos) & Enunciado \S3.1\\
D-02 & Docker en vez de Podman. Ning\'un contenedor corre como root & Enunciado \S3.2\\
D-03 & Secret Manager en vez de OpenBao. Credenciales est\'aticas y llaves en el proceso & Enunciado \S3.1\\
D-04 & Cloud Run termina el TLS. Se acepta TLS 1.2; no hay l\'imite de tasa general & Enunciado \S3.1\\
D-05 & Autenticaci\'on propia en vez de Keycloak. Sin llaves de acceso: TOTP no resiste el \emph{phishing} & Decisi\'on del equipo\\
D-06 & OPA como contenedor adjunto. F4 pasa de frontera de red a frontera de proceso & Cloud Run\\
D-07 & Continuidad reescrita. Diez a\~nos solo para la bit\'acora; retenci\'on sin bloquear & Plan gratuito de Neon\\
D-08 & Sin restricci\'on por origen. La credencial es la \'unica barrera de red & Plan gratuito de Neon\\
D-09 & La verificaci\'on de un certificado no pide sesi\'on. No revela nada ni toca la base & Enunciado \S2.2\\
D-10 & Firma Ed25519, que P1 no dise\~n\'o. La clave p\'ublica en el documento se compara con la huella publicada & Enunciado \S2.2\\
D-11 & Las tareas peri\'odicas en un solo servicio. Una falla las detiene todas & Capa gratuita\\
D-12 & Menos herramientas de an\'alisis. Sin an\'alisis est\'atico ni din\'amico & Tiempo\\
D-13 & Segundo factor obligatorio solo para DG, DC, PA y VF & Decisi\'on del equipo\\
D-14 & El aviso de bloqueo (RC-01) va a la bit\'acora y al siguiente ingreso, no por correo & Alcance\\
D-15 & Las solicitudes de cambio de cuenta tienen due\~no fuera de la matriz; la base impide autoaprobar & Hueco del dise\~no\\
D-16 & Las cuentas nacen con contrase\~na temporal y cambio obligatorio & Hueco del dise\~no\\
D-17 & Un per\'iodo se crea abri\'endolo & Hueco del dise\~no\\
D-18 & El t\'ecnico solo registra mediciones en su zona. Restringe, no ampl\'ia & Refuerzo de la matriz\\
D-19 & El transmisor corre cada 30 minutos. Un evento tarda hasta media hora en sellarse & Plan gratuito de Neon\\
D-20 & C9 se ata a la finca y C6 compara interiores de pol\'igonos & Hueco del dise\~no\\
D-21 & Anular un certificado y autorizar a un tercero piden reautenticaci\'on & Refuerzo de la matriz\\
D-22 & Un solo per\'iodo abierto a la vez, impuesto en el servidor & Hueco del dise\~no\\
D-23 & Plantillas del servidor en vez de Angular. CSP estricta, sin npm & Enunciado \S2.1\\
D-24 & El reporte semanal de HBAC-02 queda fuera; los datos est\'an en la bit\'acora & Alcance\\
D-25 & Las exportaciones del auditor son CSV en la misma aplicaci\'on, con C7 & Alcance\\
\end{xltabular}
}

\paragraph{Las de m\'as peso.} Tres desviaciones bajan una garant\'ia que la
Parte~1 promet\'ia, y conviene decirlo sin rodeos. \textbf{D-03}: sin OpenBao,
la contrase\~na de la base es est\'atica (si se filtra, sirve hasta que se
rote) y las llaves del HMAC, del cifrado y de la firma est\'an en la memoria del
proceso que las usa; quien controle el transmisor podr\'ia reescribir la cadena
de la bit\'acora. Lo acotan la instancia separada donde solo se inserta, la
comprobaci\'on de la cadena y la copia semanal que ya no se puede tocar.
\textbf{D-05}: la autenticaci\'on es c\'odigo nuestro y se pierden las llaves
de acceso, que eran la defensa resistente al \emph{phishing} contra AE-02; la
reemplazan TOTP en los cuatro roles que aprueban, emiten o cambian permisos,
la consulta a filtraciones de CO-02 y el bloqueo de RC-01. \textbf{D-07}: la
base principal solo se puede restaurar seis horas atr\'as; la retenci\'on de
diez a\~nos (BI-04, LG-06) se cumple para la bit\'acora, que guarda toda la
historia de qu\'e pas\'o con cada certificado, y con una retenci\'on que un
administrador del proyecto todav\'ia podr\'ia quitar (RR-02).

% ============================================================================
% 6. COBERTURA DE LA MATRIZ                             Responsable: Ludwin
% Presupuesto: 2 paginas. Fuente: docs/Informe_P2.md, seccion 6.
%
% Pide el enunciado: "Tabla que muestre, permiso por permiso, donde esta
% implementado y como se comprobo." Y en la evaluacion: "Permisos de la
% matriz sin implementar: se descuenta proporcionalmente".
%
% La tabla de Informe_P2.md ya esta completa (O1 a O11). Las filas de O5 y
% O7 ya estan escritas abajo. Si no cabe en dos paginas: \footnotesize, o
% quitar la columna de prueba y poner una sola linea por objeto.
% ============================================================================
\section{Cobertura de la matriz}
\label{sec:cobertura}

% Borrador pasado desde docs/Informe_P2.md (seccion 6). Ludwin: revisar los
% numeros contra tests/matriz y recortar si no cabe.
La matriz tiene 88 celdas (8 roles $\times$ 11 objetos). Entre sus pares de rol
y operaci\'on hay 84 permisos concedidos, 37 de ellos con una de las
condiciones C1 a C11, y 884 pares negados por omisi\'on. \textbf{Todos} est\'an
probados en el motor por \ruta{tests/matriz/test_celdas.py}: 977 pruebas
contra un OPA real en cada \emph{pull request}. La \cref{tab:cobertura} dice
adem\'as d\'onde vive cada permiso y qu\'e prueba lo ejercita de punta a punta,
contra PostgreSQL y OPA.

{\footnotesize
\begin{xltabular}{\linewidth}{@{}P{1.9cm}P{0.6cm}P{3.4cm}YP{3.0cm}@{}}
\caption{Cobertura de la matriz, permiso por permiso. Entre par\'entesis, la
condici\'on de la Parte~1.}
\label{tab:cobertura}\\
\toprule
\thd{Objeto} & \thd{Op.} & \thd{Roles (condici\'on)} & \thd{Endpoint} & \thd{Prueba}\\
\midrule
\endfirsthead
\contprev{5}
\toprule
\thd{Objeto} & \thd{Op.} & \thd{Roles (condici\'on)} & \thd{Endpoint} & \thd{Prueba}\\
\midrule
\endhead
\midrule
\contnext{5}
\endfoot
\bottomrule
\endlastfoot
O1 Fincas & L & DG, DC, PA; TC, VF (C1); AE (C7); DF (C8) & \ruta{GET /fincas}, \ruta{/fincas/{id}} & \ruta{test_fincas.py}\\
 & C, M & PA & \ruta{POST /fincas}, \ruta{PATCH /fincas/{id}} & \textit{idem}\\
\addlinespace
O2 Per\'iodos & L & DG, DC, TC, VF, DF; AE (C7) & \ruta{GET /periodos}, \ruta{/periodos/{id}} & \ruta{test_periodos.py}\\
 & Ab & DC & \ruta{POST /periodos}, \ruta{/periodos/{id}/abrir} (D-17, D-22) & \textit{idem}\\
 & Cr & DC & \ruta{POST /periodos/{id}/cerrar} & \textit{idem}\\
\addlinespace
O3 Mediciones & L & DG, DC, VF; TC (C1); AE (C7); DF (C8) & \ruta{GET /mediciones}, \ruta{/mediciones/{id}} & \ruta{test_mediciones.py}\\
 & C & TC (C2, y C1 por D-18); VF (C1 $\wedge$ C2) & \ruta{POST /mediciones} & \textit{idem}\\
 & A, D & VF (C3) & \ruta{POST /mediciones/{id}/aprobar}, \ruta{/devolver} & \textit{idem}\\
 & X & AE (C7) & \ruta{GET /mediciones/exportar} & \textit{idem}\\
\addlinespace
O4 Enmiendas & L & DG, DC, VF; TC (C1); AE (C7); DF (C8) & \ruta{GET /enmiendas}, \ruta{/enmiendas/{id}} & \ruta{test_enmiendas.py}\\
 & C & TC (C11) & \ruta{POST /enmiendas} & \textit{idem}\\
 & A, D & VF (C4) & \ruta{POST /enmiendas/{id}/aprobar}, \ruta{/devolver} & \textit{idem}\\
 & Cr & DC (C5) & \ruta{POST /enmiendas/{id}/cerrar} & \textit{idem}\\
 & X & AE (C7) & \ruta{GET /enmiendas/exportar} & \textit{idem}\\
\addlinespace
O5 Certificados & L & DG, DC, VF; AE (C7); DF (C8); TA (C9, atada a la finca, D-20) & \ruta{GET /certificados}, \ruta{/certificados/{id}} & \ruta{test_certificados.py}\\
 & Em & VF (C6) & \ruta{POST /certificados/{medicion}/emitir} & \textit{idem}\\
 & An & DC & \ruta{POST /certificados/{id}/anular} & \textit{idem}\\
 & X & AE (C7) & \ruta{GET /certificados/exportar} & \textit{idem}\\
\addlinespace
O6 Datos personales & L & DG, PA; DF (C8) & \ruta{GET /personas}, \ruta{/personas/{id}} & \ruta{test_flujos.py}\\
 & C, M & PA & \ruta{POST /personas}, \ruta{PATCH /personas/{id}} & \textit{idem}\\
\addlinespace
O7 Consenti\-mientos & L & DG, PA; DF (C8) & \ruta{GET /consentimientos}, \ruta{/{id}} & \ruta{test_consentimientos.py}\\
 & C & PA; DF (C8) & \ruta{POST /consentimientos} & \textit{idem}\\
 & M & DF (C8) & \ruta{PATCH /consentimientos/{id}} & \textit{idem}\\
\addlinespace
O8 Asignaciones & L & DG, DC; TC, VF (C1) & \ruta{GET /asignaciones}, \ruta{/asignaciones/{id}} & \ruta{test_asignaciones.py}\\
 & C, M & DC & \ruta{POST /asignaciones}, \ruta{PATCH /asignaciones/{id}} & \textit{idem}\\
\addlinespace
O9 Cuentas & L & DG; PA (C10) & \ruta{GET /cuentas} & \ruta{test_flujos.py}\\
 & C, M & PA (C10) & \ruta{POST /cuentas}, \ruta{PATCH /cuentas/{id}} y las solicitudes (D-15) & \textit{idem}\\
\addlinespace
O10 Credencial propia & M & los ocho roles & \ruta{PATCH /credenciales} & \ruta{test_flujos.py}\\
\addlinespace
O11 Bit\'acora & L & DG, DC, PA; AE (C7) & \ruta{GET /bitacora} & \ruta{test_bitacora.py}\\
 & X & DG, DC; AE (C7) & \ruta{GET /bitacora/exportar} & \textit{idem}\\
\end{xltabular}
}

Nadie tiene E (eliminar) sobre ning\'un objeto, y no existe ninguna ruta
DELETE. O10 no tiene L a prop\'osito: las credenciales se cambian y no se
consultan, as\'i que no existe \ruta{GET /credenciales}. Adem\'as de las
pruebas, cada objeto se ejercit\'o a mano en el sistema desplegado con las
cuentas de prueba, y los intentos negados aparecen en la propia bit\'acora
(\cref{fig:ev-denegado}).

% ============================================================================
% 7. EVIDENCIA                                          Responsable: Meibel
% Pide el enunciado: "Capturas de: el sistema funcionando por HTTPS, la
% configuracion de Secret Manager, Dependabot activo, un certificado
% verificado correctamente y el intento fallido de un rol al hacer algo que
% no le corresponde." Son las cinco de abajo (E1 a E5).
%
% Imagenes en img/: ev1-https.png, ev2-secret-manager.png, ev3-dependabot.png,
% ev4-certificado-valido.png, ev5-acceso-denegado.png
% Si el informe pasa de 10 paginas, bajar \anchoevidencia a 0.75\linewidth.
% ============================================================================
\section{Evidencia}
\label{sec:evidencia}

Las capturas se tomaron sobre el sistema desplegado en \url{\urlSistema}, desde una red distinta a la del equipo de desarrollo.

% Ancho de cada captura sola, y de cada una de un par
\newcommand{\anchoevidencia}{0.5\linewidth}
\newcommand{\anchopar}{0.29\linewidth}

% \evidencia{archivo}{etiqueta}{titulo}{descripcion}
% Titulo, descripcion e imagen van juntos, sin partirse entre paginas.
\newcommand{\evidencia}[4]{%
  \par\addvspace{0.6em}
  \noindent\begin{minipage}{\linewidth}
    {\large\bfseries\color{tecnavy} #3\par}
    \smallskip
    #4\par
    \medskip
    \centering
    \safeimage[width=\anchoevidencia]{#1}
    \captionof{figure}{#3}
    \label{#2}
  \end{minipage}\par}

% \evidenciapar{archivo1}{etiqueta1}{titulo1}{descripcion1}{archivo2}{etiqueta2}{titulo2}{descripcion2}
% Dos evidencias distintas lado a lado, cada una con su titulo, su texto, su
% captura y su propio numero de figura (asi \cref sigue apuntando a cada una).
\newcommand{\evidenciapar}[8]{%
  \par\addvspace{0.6em}
  \noindent\begin{minipage}[t]{0.485\linewidth}
    {\normalsize\bfseries\color{tecnavy} #3\par}
    \smallskip
    {\small #4\par}
    \medskip
    \centering
    \safeimage[width=\linewidth]{#1}
    \captionof{figure}{#3}
    \label{#2}
  \end{minipage}\hfill
  \begin{minipage}[t]{0.485\linewidth}
    {\normalsize\bfseries\color{tecnavy} #7\par}
    \smallskip
    {\small #8\par}
    \medskip
    \centering
    \safeimage[width=\linewidth]{#5}
    \captionof{figure}{#7}
    \label{#6}
  \end{minipage}\par}

% \evidenciados{etiqueta}{titulo}{descripcion}{archivo1}{pie1}{archivo2}{pie2}
% Una evidencia con dos capturas lado a lado, (a) a la izquierda y (b) a la
% derecha.
\newcommand{\evidenciados}[7]{%
  \par\addvspace{0.6em}
  \noindent\begin{minipage}{\linewidth}
    \captionsetup{type=figure}% para que las subfiguras funcionen aqui
    {\large\bfseries\color{tecnavy} #2\par}
    \smallskip
    #3\par
    \medskip
    \centering
    \begin{subfigure}[t]{\anchopar}
      \centering
      \safeimage[width=\linewidth]{#4}
      \caption{#5}
      \label{#1-a}
    \end{subfigure}\hspace{0.02\linewidth}
    \begin{subfigure}[t]{\anchopar}
      \centering
      \safeimage[width=\linewidth]{#6}
      \caption{#7}
      \label{#1-b}
    \end{subfigure}
    \caption{#2}
    \label{#1}
  \end{minipage}\par}


% \evidenciatres{etiqueta}{titulo}{descripcion}{archivo1}{pie1}{archivo2}{pie2}{archivo3}{pie3}
% Tres capturas: (a) y (b) lado a lado, y (c) debajo, a lo ancho (es la
% bitacora, que es una captura ancha).
\newcommand{\evidenciatres}[9]{%
  \par\addvspace{0.6em}
  \noindent\begin{minipage}{\linewidth}
    \captionsetup{type=figure}% para que las subfiguras funcionen aqui
    {\large\bfseries\color{tecnavy} #2\par}
    \smallskip
    #3\par
    \medskip
    \centering
    \begin{subfigure}[t]{\anchopar}
      \centering
      \safeimage[width=\linewidth]{#4}
      \caption{#5}
      \label{#1-a}
    \end{subfigure}\hspace{0.02\linewidth}
    \begin{subfigure}[t]{\anchopar}
      \centering
      \safeimage[width=\linewidth]{#6}
      \caption{#7}
      \label{#1-b}
    \end{subfigure}\hspace{0.02\linewidth}
    \begin{subfigure}[t]{0.38\linewidth}
      \centering
      \safeimage[width=\linewidth]{#8}
      \caption{#9}
      \label{#1-c}
    \end{subfigure}
    \caption{#2}
    \label{#1}
  \end{minipage}\par}


\evidenciapar
{ev1-https.png}{fig:ev-https}{E1. El sistema funcionando por HTTPS}{%
El sistema abre desde internet en su direcci\'on p\'ublica de Cloud Run. El
candado y el detalle del certificado muestran que la conexi\'on va cifrada
con TLS, con un certificado emitido para el dominio del servicio. Es el
\'unico servicio expuesto: el backend y el transmisor no tienen URL abierta.}
{ev2-secret-manager.png}{fig:ev-secretos}{E2. La configuraci\'on de Secret Manager}{%
Los secretos del sistema (cadenas de conexi\'on a Neon, llaves de cifrado y
de firma) viven en Secret Manager del proyecto \code{verdeazul-g4}, no en el
c\'odigo ni en el repositorio. Cada servicio los recibe al arrancar y solo
lee los suyos. La captura muestra la lista sin abrir ning\'un valor.}

% E3 y E4 en un mismo bloque: E3 a la izquierda con sus dos capturas una
% sobre la otra, E4 a la derecha. Cada una conserva su numero de figura.
\par\addvspace{0.6em}
\noindent\begin{minipage}[t]{0.485\linewidth}
  \captionsetup{type=figure}
  {\normalsize\bfseries\color{tecnavy} E3. Dependabot activo\par}
  \smallskip
  {\small Dependabot revisa las dependencias de Python, las im\'agenes de
  Docker y las acciones de GitHub, y abre un \emph{pull request} por cada
  versi\'on nueva o vulnerabilidad conocida, que pasa por el mismo CI. Ya
  encontr\'o una real, CVE-2025-62727 en Starlette, que se atendi\'o y se
  document\'o en \ruta{docs/Dependabot.md}.\par}
  \medskip
  \begin{subfigure}{\linewidth}
    \centering
    \safeimage[width=\linewidth]{ev3-dependabot-alertas.jpg}
    \caption{Alertas de Dependabot activas}
    \label{fig:ev-dependabot-a}
  \end{subfigure}

  \smallskip
  \begin{subfigure}{\linewidth}
    \centering
    \safeimage[width=0.8\linewidth]{ev3-dependabot-prs.png}
    \caption{\emph{Pull requests} abiertos por Dependabot}
    \label{fig:ev-dependabot-b}
  \end{subfigure}
  \caption{E3. Dependabot activo}
  \label{fig:ev-dependabot}
\end{minipage}\hfill
\begin{minipage}[t]{0.485\linewidth}
  {\normalsize\bfseries\color{tecnavy} E4. Un certificado verificado correctamente\par}
  \smallskip
  {\small El verificador \code{vf2} emiti\'o y firm\'o el certificado de la
  medici\'on~1, \code{df1} (el due\~no de la finca) lo descarg\'o, y la
  p\'agina p\'ublica de verificaci\'on lo comprob\'o sin iniciar sesi\'on. La
  firma Ed25519 coincide con el contenido, y la huella de la llave que firm\'o
  es la publicada por Verdeazul.\par}
  \medskip
  \centering
  \safeimage[width=\linewidth]{ev4-certificado-valido.png}
  \captionof{figure}{E4. Un certificado verificado correctamente}
  \label{fig:ev-valido}
\end{minipage}\par

\evidenciatres{fig:ev-denegado}{E5. Intento fallido de un rol}{%
Cada rol solo ve en pantalla lo que la matriz le concede. El técnico de campo tc1 ingresa y escribe a mano una direcci\'on que no le corresponde (certificados), el servidor consulta el motor de pol\'iticas, responde \enquote{Acceso denegado} y registra el intento en la bit\'acora.}
{ev5-panel-tc1.png}{Panel del técnico 1}
{ev5-acceso-denegado.png}{Acceso denegado}
{ev5-bitacora.png}{Bitácora}

% ============================================================================
% ANEXO A. REGISTRO DE DECISIONES DE ARQUITECTURA      Responsable: Adriel
% Reemplaza a docs/P2_Desviaciones_respecto_P1.md (ADR de Ludwin, revisados
% contra la Parte 1). Queda fuera del conteo de 6 a 10 paginas: es el detalle
% que la seccion 5 resume.
% ============================================================================
\clearpage
\appendix
\crefname{section}{anexo}{anexos}
\Crefname{section}{Anexo}{Anexos}
\section{Registro de decisiones de arquitectura}
\label{anx:adr}

El diseño de octubre asumía servidores propios y una restricción de la
cliente: solo software libre. La Parte~2 impone una plataforma concreta
(Cloud Run, Neon, Secret Manager y Artifact Registry, con imágenes de Docker),
y casi todas las desviaciones salen de ese cambio de supuesto. Cada una va
como registro de decisión de arquitectura (ADR): el contexto, lo que se
decidió y sus consecuencias. El identificador D-XX es el que usan el código
y los demás documentos del repositorio.

% \adr{numero}{titulo}: encabezado de un ADR, fuera del indice.
\newcommand{\adr}[2]{\subsection*{ADR #1. #2}}
\newcommand{\adrctx}{\textbf{Contexto.}\ }
\newcommand{\adrdec}{\par\smallskip\textbf{Decisión.}\ }
\newcommand{\adrpos}{\par\smallskip\textbf{Consecuencias positivas.}}
\newcommand{\adrneg}{\par\smallskip\textbf{Consecuencias negativas.}}

\adr{001 (D-01)}{La plataforma de alojamiento es software propietario}
\adrctx El Informe de políticas recoge una restricción explícita de la cliente:
presupuesto solo para software libre (requerimiento R11, regla LG-08). La
Parte~2 impone Cloud Run, Neon, Secret Manager y Artifact Registry, que son
servicios propietarios.
\adrdec El sistema se aloja en esos servicios. Todo el software que escribimos y
sus dependencias siguen siendo libres (Python, FastAPI, PostgreSQL, PostGIS,
OPA); lo propietario es solo dónde se ejecuta.
\adrpos
\begin{itemize}
  \item Los cuatro servicios tienen capa gratuita suficiente: se respeta la
    restricción de presupuesto, aunque no la de licencia.
  \item Siguen aplicando las mitigaciones de P1 para el riesgo del proveedor:
    los datos personales se cifran en la aplicación y la base guarda texto
    cifrado.
\end{itemize}
\adrneg
\begin{itemize}
  \item Crece RR-02 de P1 (el proveedor o una cuenta de administración de la
    nube comprometida accede a los datos): ahora el proveedor también custodia
    los secretos y ejecuta el motor de base de datos.
  \item R11 y LG-08 no se cumplen en la capa de alojamiento.
\end{itemize}

\adr{002 (D-02)}{Docker en lugar de Podman sin privilegios}
\adrctx La arquitectura eligió Podman y descartó Docker Engine porque su
servicio central corre como administrador. La Parte~2 exige imágenes de
Docker en Artifact Registry.
\adrdec Las imágenes se construyen con Docker. Cada Dockerfile declara un
usuario sin privilegios (\code{USER 10001:10001} en api, edge y transmisor;
\code{USER 1000:1000} en el motor de políticas), así que ningún contenedor
corre como root.
\adrpos
\begin{itemize}
  \item El demonio de Docker existe solo donde se construye la imagen, en los
    \emph{runners} de GitHub Actions. En producción ejecuta Cloud Run, que no
    usa Docker y aísla cada instancia.
  \item Se conserva el mínimo privilegio en la ejecución.
\end{itemize}
\adrneg
\begin{itemize}
  \item La construcción depende de un demonio con privilegios de
    administrador, aunque sea en una máquina efímera de la integración
    continua.
\end{itemize}

\adr{003 (D-03)}{Secret Manager en lugar de OpenBao}
\adrctx P1 había descartado los gestores de secretos de los proveedores de nube
por crear dependencia, y la Parte~2 los impone. OpenBao hacía tres cosas que
Secret Manager no hace: generaba credenciales de base de datos de vida
limitada, firmaba cada eslabón de la bitácora con su motor Transit sin que la
llave saliera de él, y cifraba los datos personales sin entregarle la llave a
nadie.
\adrdec Los secretos viven en Secret Manager y se inyectan como variables de
entorno; cada servicio tiene su propia cuenta y solo lee lo suyo. Las tres
funciones de OpenBao pasan a la aplicación: una contraseña de Neon guardada
como secreto, con rotación manual; el HMAC de la bitácora calculado en el
proceso del transmisor; y el cifrado de cédula y teléfono (LG-05) en el
proceso del backend, con la búsqueda por cédula sobre un valor derivado.
\adrpos
\begin{itemize}
  \item Ninguna credencial vive en el repositorio ni en las imágenes, y cada
    servicio solo lee lo suyo.
  \item Los cambios de configuración y de permisos sobre los secretos quedan
    en los registros de auditoría de administración de Google Cloud, activos
    por defecto.
  \item El rol de la aplicación no es dueño de las tablas ni puede alterar el
    esquema, y la bitácora vive en otra instancia donde el transmisor solo
    inserta.
  \item La tarea \ruta{volcado} (ADR 007) sube cada semana la bitácora completa,
    con el eslabón de cada evento, a un depósito donde nada se puede borrar ni
    reemplazar en diez años. Si alguien con la llave reescribiera la cadena en
    Neon, dejaría de coincidir con las copias anteriores, que ya no puede
    tocar. La tarea \ruta{cadena} la recalcula y deja constancia en la propia
    bitácora.
  \item Desaparece RR-05 de P1 (el nodo único de OpenBao): ese riesgo pasa a
    ser parte de RR-02.
\end{itemize}
\adrneg
\begin{itemize}
  \item Si la credencial de la base se filtra, sirve hasta que alguien la
    cambie. Con OpenBao caducaba en 24 horas.
  \item Quien comprometa el proceso del transmisor puede reescribir la cadena
    en Neon. La garantía baja de «nadie con acceso a la base puede falsificar
    el historial» a «nadie con acceso a la base, pero sí quien controle el
    transmisor»; las copias semanales anteriores siguen intactas.
  \item Las llaves de los datos personales y de la firma están en la memoria
    del proceso, no encerradas en un motor que solo responde operaciones.
  \item Las \emph{lecturas} de un secreto no quedan registradas: el registro
    de acceso a datos de Google Cloud viene apagado y no se activó.
  \item El sistema queda atado al proveedor, que era lo que P1 quería evitar.
\end{itemize}

\adr{004 (D-04)}{Cloud Run termina el TLS}
\adrctx La arquitectura ponía nginx como única entrada, a cargo de terminar
TLS~1.3, limitar tamaño y tasa y agregar las cabeceras de seguridad. Cloud
Run asigna un dominio con TLS automático y termina la conexión antes del
contenedor.
\adrdec F1 la aplica Google. Las cabeceras (CSP, \code{X-Frame-Options},
\code{X-Content-Type-Options}, \code{Referrer-Policy},
\code{Permissions-Policy} y \code{Strict-Transport-Security}) y los límites
de tamaño pasan al servicio de borde. F2 se implementa con los permisos de
Cloud Run: el backend no acepta invocaciones anónimas y solo la cuenta del
borde puede llamarlo.
\adrpos
\begin{itemize}
  \item No hay certificados que gestionar ni renovar.
  \item F2 es más fuerte que la red privada diseñada: quien conozca la URL del
    backend recibe un rechazo antes de que se ejecute nuestro código.
\end{itemize}
\adrneg
\begin{itemize}
  \item La versión de TLS y los cifrados quedan fuera de nuestro control:
    Cloud Run acepta TLS~1.2 y 1.3, y P1 exigía solo 1.3.
  \item No se reimplementó el límite de tasa general de nginx. Solo la
    verificación pública tiene uno (ADR 009), y el inicio de sesión queda
    acotado por el bloqueo de RC-01.
\end{itemize}

\adr{005 (D-05)}{Autenticación propia en lugar de Keycloak}
\adrctx P1 eligió Keycloak como único emisor de identidades y advirtió, con
razón, que escribir la autenticación uno mismo es un error frecuente. Pero
Keycloak necesita su propia base, arranca lento, consume memoria y expone una
consola de administración: desproporcionado para un servicio que escala a
cero en una capa gratuita.
\adrdec La autenticación pasa al backend, con Argon2id y los parámetros de
P1 (variante id, sal por usuario, 19~MiB, dos iteraciones, un hilo). Se
conserva el diseño de sesión de P1: cookie opaca \code{HttpOnly},
\code{Secure} y \code{SameSite}, sesión en la base, revocación inmediata al
dar de baja o cambiar de rol, e identificador nuevo en cada ingreso. CO-01 a
CO-05 y el bloqueo de RC-01 viven en la aplicación. El segundo factor es TOTP
(ADR 013), con reautenticación en las acciones críticas.
\adrpos
\begin{itemize}
  \item Un componente menos que operar y ninguna consola de administración
    expuesta.
  \item Se conserva lo valioso del diseño, que era la sesión.
\end{itemize}
\adrneg
\begin{itemize}
  \item Se pierden las llaves de acceso, que P1 exigía para las direcciones,
    el personal administrativo y los verificadores; implementar WebAuthn bien
    no cabía en el tiempo. Con ellas se pierde la resistencia al
    \emph{phishing}: un código TOTP se puede capturar con un sitio falso en
    tiempo real.
  \item AE-02 (contraseñas filtradas o robadas) queda cubierta por TOTP en los
    cuatro roles que aprueban, emiten o cambian permisos, por CO-02 (lista de
    prohibidas y consulta a filtraciones conocidas) y por el bloqueo de RC-01,
    pero no por un factor resistente al \emph{phishing}.
  \item El código de autenticación es nuestro, que es el riesgo que P1
    advertía. Lo acotan sus pruebas en \ruta{tests/identidad}.
\end{itemize}

\adr{006 (D-06)}{OPA como contenedor adjunto}
\adrctx El motor de políticas era un componente aparte que solo aceptaba
consultas desde la red privada (F4). Cloud Run no ofrece red privada entre
servicios sin configuración de red que la capa gratuita no cubre bien.
\adrdec OPA corre como contenedor adjunto dentro del mismo servicio que el
backend, escuchando solo en la interfaz local. Las políticas van dentro de su
imagen y llevan la misma etiqueta de versión.
\adrpos
\begin{itemize}
  \item Se conserva la separación entre el punto de aplicación y el de
    decisión: las políticas siguen siendo código en Rego, versionado y con
    pruebas propias, y el backend no puede conceder lo que el motor niega.
  \item Se sabe qué reglas están corriendo, por la etiqueta de la imagen.
\end{itemize}
\adrneg
\begin{itemize}
  \item F4 deja de ser frontera de red y pasa a ser de proceso: quien
    comprometa el contenedor del backend está en la misma instancia que el
    motor.
\end{itemize}

\adr{007 (D-07)}{Las reglas de continuidad se reescriben}
\adrctx DI-01 a DI-04 prometían recuperación en dos horas, respaldos cada hora,
regla 3-2-1 con una copia aislada imborrable y pérdida máxima de una hora,
para dos máquinas propias con réplica en espera. Neon en plan gratuito
restaura a cualquier punto de las últimas seis horas, y las instantáneas
programadas son de pago. Seis horas no sirven para BI-04 ni para LG-06, que
exigen diez años para lo ligado a un certificado.
\adrdec DI-01 y DI-04 se cubren con la restauración de Neon. La retención la
resuelve la tarea \ruta{volcado}, que cada lunes copia la \textbf{bitácora}
a un depósito de Cloud Storage con retención de diez años: nada se borra ni
se reemplaza antes, y la cuenta del transmisor solo puede crear objetos. La
copia lleva el eslabón de cada evento y un manifiesto con su SHA-256.
\adrpos
\begin{itemize}
  \item La restauración de Neon cumple DI-04 con margen y DI-01, porque
    restaurar es inmediato.
  \item El volcado recupera la copia separada de la bitácora de BI-03, su
    retención de diez años (BI-04, LG-06) y, para la bitácora, la copia fuera
    de la cuenta de Neon (DI-02) que no se puede borrar (DI-03). Se verifica
    sin acceso a Neon, con el archivo y la llave.
\end{itemize}
\adrneg
\begin{itemize}
  \item \textbf{Solo la bitácora.} Copiar la base principal exigía un rol con
    lectura de todo el esquema, incluidos los datos personales cifrados. Lo
    que exige diez años es la trazabilidad, y esa está entera en la bitácora.
    Si la base principal se perdiera más allá de seis horas, los datos
    vigentes habría que reconstruirlos, pero no se pierde la historia de qué
    pasó con ellos. DI-02 y DI-03 se cumplen para la bitácora y no para la
    base principal.
  \item \textbf{La retención no está bloqueada.} Bloquearla es irreversible:
    el depósito no se podría borrar en diez años y Google le pondría un
    gravamen al proyecto. Un administrador del proyecto podría quitar la
    política y borrar las copias (RR-02). En producción se cierra con
    \code{-{}-lock-retention-period}.
  \item Se pierden la réplica en espera y el ensayo semestral de recuperación;
    la disponibilidad deja de estar bajo nuestro control, y crece RR-04.
\end{itemize}

\adr{008 (D-08)}{Sin restricción por dirección de origen}
\adrctx La lista de direcciones permitidas y el acceso por red privada son de
plan pago en Neon. Además, Cloud Run no sale a internet por una dirección
estable sin configurar salida a una red virtual y traducción de direcciones.
\adrdec La base acepta conexiones desde cualquier lugar, y la barrera es la
credencial.
\adrpos
\begin{itemize}
  \item Neon exige TLS y autentica con SCRAM, y la aplicación exige además
    \code{sslmode=verify-full}: la credencial no se puede capturar en tránsito
    y se verifica que el servidor sea el de Neon.
\end{itemize}
\adrneg
\begin{itemize}
  \item No hay segunda línea de defensa en la red: todo recae en que la
    credencial no se filtre. Eso sube la importancia de que nada se versione
    (gitleaks en cada pull request) y de que cada rol tenga el mínimo
    privilegio.
\end{itemize}

\adr{009 (D-09)}{La verificación de un certificado no exige sesión}
\adrctx El diseño sostenía la mediación completa: nada salta el punto de
aplicación. El enunciado pide que cualquiera pueda comprobar un certificado
sin acceso al sistema, porque un certificado circula fuera de Verdeazul y
vale dinero.
\adrdec La verificación es pública. No consulta la base: recibe el archivo,
comprueba la firma con la clave pública que trae y compara su huella con la
de Verdeazul (ADR 010). La acotan un tamaño máximo de 64~KB (en el borde y en
el backend), un límite de 10 verificaciones por minuto por origen en el
borde, y una respuesta que solo dice si es válido y las dos huellas. La
huella de la llave de producción es
\ruta{SHA256:203b7265f5b57cbe0aed5ef62cb74a56e2731c5185675e8742203dfbd3df811d}.
\adrpos
\begin{itemize}
  \item No hay objeto protegido en juego: la verificación no revela nada.
  \item Un tercero puede confiar sin acceso al sistema. Comprobado en
    producción el 17 de noviembre: un certificado firmado con otra llave tiene
    firma válida y aun así se rechaza, porque su huella es otra.
\end{itemize}
\adrneg
\begin{itemize}
  \item El límite de tasa vive en la memoria de cada instancia del borde; con
    varias, se multiplica. En una operación real iría en un almacén compartido.
  \item Una verificación válida no dice si el certificado se anuló después,
    porque un anulado conserva su firma. Eso se le pregunta a Verdeazul.
\end{itemize}

\adr{010 (D-10)}{Firma digital de certificados}
\adrctx P1 no lo diseñó: la Ley 8454 aparecía en el marco regulatorio, pero sin
formato, llave ni procedimiento de verificación. El enunciado lo exige.
\adrdec Cada certificado se emite firmado con Ed25519. La clave privada es un
secreto de Secret Manager que solo lee el backend desplegado; la pública viaja
en el documento. El evento de emisión en la bitácora lleva la huella de la
clave, un campo nuevo respecto a la Tabla 4.2 de P1.
\adrpos
\begin{itemize}
  \item Cualquiera comprueba que el contenido no cambió desde que se firmó.
  \item Nadie del equipo necesita la llave de producción: en local se firma
    con una llave de prueba.
\end{itemize}
\adrneg
\begin{itemize}
  \item La llave privada está en la memoria del backend (ADR 003), no en un
    módulo que solo firme.
  \item La clave pública dentro del documento prueba que el contenido no
    cambió, pero no quién firmó: cualquiera puede generar un par y firmar un
    certificado falso que verifica igual. Se cierra publicando la huella de la
    clave de Verdeazul por un canal independiente. La verificación pública la
    compara sola; quien use herramientas propias, a mano.
\end{itemize}

\adr{011 (D-11)}{Los trabajos programados se consolidan en un solo servicio}
\adrctx P1 preveía varios procesos periódicos: el transmisor, la desactivación
de lo vencido (HBAC-03), el reporte semanal de HBAC-02 y la comprobación
mensual de la cadena (BI-02). Cloud Run cobra por instancia activa, y la capa
gratuita de Cloud Scheduler cubre tres trabajos.
\adrdec Todas las tareas viven en el transmisor, con un punto de entrada cada
una (\ruta{POST /tareas/{nombre}}), que exige la cuenta del planificador y no
es accesible desde internet. Hay tres trabajos programados:
\ruta{transmitir} cada 30 minutos (ADR 019), \ruta{volcado} los lunes y
\ruta{cadena} el primer día de cada mes, que es la comprobación de BI-02.
\ruta{vencer} se lanza a mano.
\adrpos
\begin{itemize}
  \item Un solo servicio y tres trabajos, dentro de la capa gratuita.
  \item La bandeja de salida tolera atrasos: los eventos se acumulan sin
    perderse.
  \item Que \ruta{vencer} no corra sola no abre ningún acceso: C1 lee en cada
    petición solo las asignaciones vigentes, y C9 compara la fecha en el motor.
\end{itemize}
\adrneg
\begin{itemize}
  \item Una falla de ese servicio detiene todas las tareas a la vez.
  \item Hasta que alguien lanza \ruta{vencer}, una asignación o autorización
    vencida no queda marcada como revocada ni deja su evento en la bitácora.
  \item El reporte de HBAC-02 no se implementó (ADR 024).
\end{itemize}

\adr{012 (D-12)}{Menos herramientas de análisis}
\adrctx Las fichas de tecnología de P1 listaban Bandit, Semgrep, pip-audit,
Trivy, OWASP ZAP, gitleaks, hadolint, Gixy, Regal y testssl.sh. En tres
semanas no se configuran todas con sentido.
\adrdec Se mantienen tres: Dependabot, que el enunciado exige y cubre lo de
pip-audit; gitleaks sobre todo el historial en cada pull request; y las
pruebas de autorización por cada celda de la matriz, que P1 ya exigía.
\adrpos
\begin{itemize}
  \item Cubren lo que el enunciado más castiga: una credencial versionada
    anula el criterio de secretos, y el historial de git conserva lo que se
    borra. Dependabot ya detectó una vulnerabilidad real, CVE-2025-62727 en
    Starlette, que se atendió y se documentó en \ruta{docs/Dependabot.md}.
  \item Las pruebas de la matriz son la evidencia de la tabla de cobertura.
\end{itemize}
\adrneg
\begin{itemize}
  \item No hay análisis estático del código ni de las imágenes, ni dinámico
    de la aplicación desplegada.
\end{itemize}

\adr{013 (D-13)}{El segundo factor es obligatorio solo para cuatro roles}
\adrctx La Tabla 7.1 del Informe de políticas pedía segundo factor a los ocho
perfiles, con llave de acceso preferida y TOTP como alternativa. El ADR 005
explica por qué no hay llaves de acceso; esta es la otra mitad.
\adrdec El segundo factor es obligatorio para Dirección General, Dirección de
Certificación, Personal Administrativo y Verificador, que son los que
aprueban, emiten o cambian permisos: sin él no pasan de la pantalla de
activación (CU-04). Para técnico de campo, auditor externo, dueño de finca y
tercero autorizado es opcional, y una vez activo se pide siempre. El conjunto
es una constante en \ruta{services/api/app/auth/totp.py}. Los códigos de
recuperación de RC-02 están implementados: ocho de un solo uso, guardados
como huella SHA-256.
\adrpos
\begin{itemize}
  \item Responde al uso real: el técnico trabaja con conexión irregular, y el
    dueño y el tercero entran pocas veces al año.
  \item Se conservan el cierre por inactividad de cada perfil y la
    reautenticación en las acciones críticas de la Tabla 7.1.
  \item Extenderlo a los ocho roles es cambiar una línea.
\end{itemize}
\adrneg
\begin{itemize}
  \item En los cuatro roles sin factor obligatorio, una contraseña filtrada
    basta para entrar (AE-02). Lo acotan CO-02, el bloqueo de RC-01 y que
    ninguno de esos roles aprueba, emite ni cambia permisos: el técnico
    registra mediciones y pide enmiendas, y una medición siempre la aprueba
    otra persona (C3).
  \item La recuperación por videollamada de RC-03 es un procedimiento humano y
    queda fuera del sistema.
\end{itemize}

\adr{014 (D-14)}{El aviso al titular de RC-01 no llega por correo}
\adrctx RC-01 dice «tras 10 intentos fallidos, bloqueo de 30 minutos y aviso al
titular». El sistema no envía correos.
\adrdec El bloqueo está implementado tal cual, con espera creciente desde el
quinto intento. El aviso toma dos formas: el evento
\ruta{CUENTA_BLOQUEADA} en la bitácora, y un mensaje en la pantalla del
siguiente ingreso con los intentos fallidos desde el último.
\adrpos
\begin{itemize}
  \item Cada bloqueo deja rastro, y agregar un correo es enviarlo desde el
    mismo punto donde hoy se registra el evento.
\end{itemize}
\adrneg
\begin{itemize}
  \item El titular se entera cuando vuelve a entrar, no en el momento.
\end{itemize}

\adr{015 (D-15)}{Las solicitudes de cambio de cuenta tienen dueño fuera de la matriz}
\adrctx C10 exige aprobación previa de Dirección General antes de que el
Personal Administrativo cree o modifique una cuenta. Esa aprobación vive en
\ruta{solicitudes_rol}, que no es un objeto de la matriz, así que la matriz
no dice quién crea ni quién aprueba una solicitud. Además, el esquema solo
modelaba el cambio de rol, y C10 gobierna también el alta y la baja.
\adrdec Con CU-01 y CU-03 de P1: el Personal Administrativo registra la
solicitud de la jefatura, Dirección General la aprueba o la rechaza, y el
Personal Administrativo la aplica. La base impide que quien pide apruebe, que
quien aprueba aplique y que nadie actúe sobre su propia cuenta. L·C10 se
implementa literal: el Personal Administrativo ve una cuenta solo mientras
tiene un cambio aprobado pendiente de aplicar. Un cambio de rol revoca las
asignaciones de zona (CU-10), y el cambio de rol y la baja cierran al
instante las sesiones (SE-01, CU-11).
\adrpos
\begin{itemize}
  \item C10 cubre alta, cambio y baja, y la separación de deberes la impone la
    base, no solo el código.
\end{itemize}
\adrneg
\begin{itemize}
  \item Quién toca las solicitudes se decide en el servidor y no en el motor:
    lo cubren las pruebas de identidad, no las de la matriz.
\end{itemize}

\adr{016 (D-16)}{Las cuentas nuevas nacen con una contraseña temporal}
\adrctx P1 no dijo cómo recibe alguien su primera contraseña; con Keycloak lo
resolvía el correo de activación.
\adrdec Al aplicar el alta, el Personal Administrativo recibe una contraseña
temporal aleatoria de 20 caracteres que se muestra una vez. La cuenta nace
con \ruta{debe_cambiar}: la sesión solo sirve para activar el segundo factor
y cambiarla.
\adrpos
\begin{itemize}
  \item El administrativo conoce la temporal, pero no la definitiva ni el
    segundo factor: no puede usar la cuenta que creó.
\end{itemize}
\adrneg
\begin{itemize}
  \item Hasta el cambio, la temporal pasó por las manos del administrativo.
\end{itemize}

\adr{017 (D-17)}{Un período se crea abriéndolo}
\adrctx La matriz le da a Dirección de Certificación Ab y Cr sobre O2, pero un
período no se puede abrir si no existe, y ningún endpoint lo creaba.
\adrdec \ruta{POST /periodos} crea un período que nace ABIERTO, bajo Ab;
\ruta{POST /periodos/{id}/abrir} reabre uno cerrado. Las dos piden
reautenticación, como cerrar. Cerrar exige que no queden mediciones
REGISTRADAS; las DEVUELTAS no cuentan, porque no se corrigen sino que se
enmiendan.
\adrpos
\begin{itemize}
  \item Ninguna celda de la matriz cambia, y cada apertura queda en la
    bitácora.
\end{itemize}
\adrneg
\begin{itemize}
  \item Reabrir un período vuelve a habilitar C2 y suspende C11 mientras está
    abierto: quien lo reabre tiene que saber si hay enmiendas en curso.
\end{itemize}

\adr{018 (D-18)}{El técnico solo registra mediciones en su zona}
\adrctx La matriz le da al técnico C sobre O3 con C2 y no menciona C1. Leída
literal, podría registrar en una finca de una zona que no tiene asignada.
\adrdec Al registrar, el servidor además le pregunta al motor si el técnico
puede leer esa finca (L sobre O1, con C1); si no, niega y lo registra.
\adrpos
\begin{itemize}
  \item Una cuenta de técnico comprometida no inyecta mediciones en fincas
    ajenas, y no se amplía ningún permiso: ABAC solo restringe.
\end{itemize}
\adrneg
\begin{itemize}
  \item La condición vive donde la matriz no la puso.
\end{itemize}

\adr{019 (D-19)}{El transmisor corre cada media hora}
\adrctx P1 describía un transmisor casi en tiempo real, con aviso si el retraso
pasaba de cinco minutos. Cada corrida despierta las dos bases de Neon por
cinco minutos: cada cinco minutos no se dormirían nunca, unas 180 horas al
mes contra las 100 gratuitas, y Neon suspendería la base.
\adrdec Cloud Scheduler invoca al transmisor cada 30 minutos. La hora que queda
en la bitácora es la de la operación, no la de la copia.
\adrpos
\begin{itemize}
  \item Las horas de cómputo alcanzan y ningún evento se pierde: la bandeja se
    escribe en la misma transacción que la operación.
\end{itemize}
\adrneg
\begin{itemize}
  \item Un evento tarda hasta media hora en sellarse (más los dos minutos de
    margen del transmisor). Para la presentación se invoca a mano.
  \item Durante esa media hora no hay detección temprana; lo acota que la
    aplicación no puede modificar la bandeja (migración 012).
\end{itemize}

\adr{020 (D-20)}{C9 se ata a la finca, y C6 compara el interior de los polígonos}
\adrctx C9 dice $\mathit{autorizacion.tercero} = \mathit{sujeto.id} \wedge
\mathit{hoy} < \mathit{autorizacion.vence}$ y no dice de qué finca: un tercero
con cualquier autorización vería los certificados de todas, lo que LG-03
prohíbe; tampoco excluye una autorización revocada (CU-09). C6 dice
«ningún certificado vigente tiene un polígono que se superponga», y con
\code{ST\_Intersects} solo, dos fincas que comparten un lindero se cruzan en
el borde.
\adrdec Para C9 el motor evalúa la condición tal cual, pero el backend solo le
manda la autorización de la finca del certificado y no revocada; la fecha la
sigue decidiendo el motor. Para C6 se exige que los interiores se crucen:
\code{ST\_Intersects} y no \code{ST\_Touches}.
\adrpos
\begin{itemize}
  \item Un tercero ve solo lo que le autorizaron, y una revocación deja de
    servir; dos fincas vecinas se pueden certificar las dos. Lo cubren
    \ruta{test_una_autorizacion_de_otra_finca_no_sirve} y
    \ruta{test_dos_fincas_vecinas_si}.
\end{itemize}
\adrneg
\begin{itemize}
  \item Una parte de C9 vive en una consulta SQL y no en Rego.
\end{itemize}

\adr{021 (D-21)}{Anular un certificado y autorizar a un tercero piden reautenticación}
\adrctx La Tabla 7.1 de P1 pide reautenticar antes de aprobar, emitir, abrir y
cerrar períodos y cambiar roles. Dos acciones más pesan igual: anular un
certificado (libera el terreno para otra emisión, así que con una sesión
robada se podría anular y reemitir) y autorizar o extender una autorización a
un tercero (entregar información del dueño, Ley 8968, art.~14).
\adrdec Esas dos acciones piden reautenticación. Revocar una autorización no,
a propósito: protege al dueño y no debe costarle nada.
\adrpos
\begin{itemize}
  \item Una sesión robada ya no alcanza para anular y reemitir ni para
    entregar datos a un tercero, y no cambia ninguna celda.
\end{itemize}
\adrneg
\begin{itemize}
  \item Un paso más en esas dos acciones.
\end{itemize}

\adr{022 (D-22)}{Un solo período de medición abierto a la vez}
\adrctx P1 define los períodos como ventanas de registro y solo exige C2. No
dice cuántos puede haber abiertos, y con dos un técnico podía registrar la
misma finca en las dos ventanas.
\adrdec Solo uno abierto: crear o reabrir responde 409 si ya hay otro, con un
candado para que dos aperturas simultáneas no dejen dos. La regla vive en el
servidor, porque un índice único obligaba a rehacer las pruebas de cinco
módulos.
\adrpos
\begin{itemize}
  \item Cada medición pertenece a una sola ventana, y no cambia ninguna celda.
\end{itemize}
\adrneg
\begin{itemize}
  \item Con la credencial de \ruta{app_rw} se podría abrir un segundo período
    por SQL.
  \item El período siguiente espera a que el anterior se cierre, y cerrar
    exige que no queden mediciones sin revisar.
\end{itemize}

\adr{023 (D-23)}{La interfaz son plantillas del servidor, no Angular}
\adrctx P1 eligió Angular. El enunciado dice que la interfaz no se evalúa y que
basta una interfaz sobria que permita ejercitar todos los permisos.
\adrdec El borde arma las páginas con plantillas Jinja, sin un marco de
JavaScript.
\adrpos
\begin{itemize}
  \item Desaparece la cadena de dependencias de npm; la CSP puede ser estricta
    (\code{script-src 'self'}); el anti-CSRF es un token por formulario
    comprobado en el borde.
\end{itemize}
\adrneg
\begin{itemize}
  \item Se pierde la usabilidad que motivó a P1: cada acción recarga la página.
\end{itemize}

\adr{024 (D-24)}{El reporte semanal de HBAC-02 queda fuera de alcance}
\adrctx HBAC-02 pide un reporte semanal de volumen anómalo. No es una celda de
la matriz, y el enunciado (§2.1) pide recortar en lo accesorio y nunca en la
matriz ni en los controles.
\adrdec No se implementa.
\adrpos
\begin{itemize}
  \item El tiempo fue a la matriz completa y a los controles obligatorios, y
    los datos para hacerlo están en la bitácora.
\end{itemize}
\adrneg
\begin{itemize}
  \item Un volumen anómalo no se detecta solo (crece RR-01). Lo acota que cada
    medición la aprueba otra persona (C3).
\end{itemize}

\adr{025 (D-25)}{Las exportaciones del auditor van en la misma aplicación, en CSV}
\adrctx La minuta comprometió un portal propio para auditores con reportes en
PDF. La matriz da X sobre mediciones, enmiendas, certificados y bitácora.
\adrdec Son cuatro rutas de la misma aplicación que devuelven CSV
(\ruta{/mediciones/exportar}, \ruta{/enmiendas/exportar},
\ruta{/certificados/exportar} y \ruta{/bitacora/exportar}), limitadas por C7
a la ventana de auditoría.
\adrpos
\begin{itemize}
  \item Las celdas X están implementadas con el mismo control de acceso que el
    resto, y hay un componente menos que proteger.
\end{itemize}
\adrneg
\begin{itemize}
  \item No hay portal separado ni PDF: el auditor procesa los datos con sus
    herramientas.
\end{itemize}

\end{document}